% !TeX root = ../papers/nc-p3.tex
\section{非正規四次の正規化の切断種数と十八次の排除}
\label{sec:cofactor-eighteen-quartic-normalization}
\begin{theorem}\label{thm:cofactor-exclude-eighteen}
被約窓で整な一般被約の階数二曲線の次数十八は生じない。
ゆえにこの整曲線の場合は次数二十の二次核曲面射影、または次数二十一の七余因子Bezout表示に限られる。
\end{theorem}
\begin{proof}[四次正規化の不正則数]
先の定理群から非正規非錐四次核曲面 $K$ だけを扱う。
その正規化の最小解消 $S$ 上で入力と核の二つの自由な束は
$H^2=8$, $HL=6$, $L^2=4$。入力切断の種数は三以下。
核切断は一般平面四次の正規化で、非正規曲面の特異曲線に会うので種数 $g_L$ は二以下である。
従って $K_S H\le-4$, $K_S L=2g_L-6\le-2$。特に $S$ は線織面と双有理である。

その底種数 $q>0$ と仮定する。直線束モデルの一般母線 $F$ と $k=HF$, $l=LF$ を使う。
任意の自由な束 $D$ の母線次数を $j>0$ とすると
\[
 K_S D\ge(2q-2)j-D^2/j .
\]
実際射影直線束の数値格子での等式に、解消時の
$\sum m_i(1-m_i/j)\ge0$ を加えればよい。従って $k,l\le2$。
$k=1$ なら入力系の全非零基点は重複度一、$k=2$ なら上式と $K_SH\le-4$ から
$q=1$ で全非零基点は重複度二である。定理\ref{thm:cofactor-eighteen-no-smooth-quadric}の初等変換で
$H$ に基点のない直線束モデルを取る。このモデルの核移動部は自乗
$2l(HL)/k-l^2 H^2/k^2$ であり、四以上である。
$k=1$ なら $l=1$ だけ、$k=2$ なら $l=1,2$ で、どれもその自乗はちょうど四である。
従って両系がこの一つの射影直線束上で自由となる。

$k=l=1$ では $H\equiv L+2F$。$L$ は豊富である。母線以外を一点に潰せば
全ての核母線が一点を通るため錐となるからである。$H$ も豊富で入力は有限。
跡の有効因子 $D_{\rm tr}\sim2L-K_S-2H$ について
$D_{\rm tr}L=2-2g_L=2-2q\le0$ だが $D_{\rm tr}F=2$、これは不可能。
$k=l=2$ では $H\equiv L+F$ は豊富である。$K_SH=-4$, $K_SL=-2$ で
$\chi(\mathcal O_S(L+tH))$ と $\chi(M(t))$ の二次・一次係数は一致し、有限評価包含は余次元一で同型となる。
深さ二から全体でも同型だが定数項は三と四で異なる。

$k=2,l=1$ では底は楕円曲線で、最小切断を $C_0$, $C_0^2=-e$ とする。
線織面の不変量の上界 $C_0^2\le q$（Segre--Nagataの切断定理）より $e\ge-1$。
$L^2=4$ から $e$ は偶数で、$LC_0=(4-e)/2$, $HC_0=2-e$。
nef性と、楕円曲線の次数一の束が基点を持つことから $e=0$。
このとき $H\equiv2C_0+2F$ は豊富であり、$g_L=1$, $K_SH=-4$。
有限評価包含は再び余次元一、従って全体で $M=f_*L$ となる。
しかし $H|_{C_0}$ は次数二の生成された束で入力直線へ二対一。
その直線上で有限繊維長は二以上となり、生成数との等号から直線が階数二軌跡に入ってしまう。
従って $S$ は有理である。
\end{proof}

\begin{proof}[切断種数一と二]
$g_L=2$ のとき、$K_S+L$ の切断空間は随伴と有理性により一般核切断の標準切断空間に同型である。
移動部 $U$ は $LU=2$（標準系は自由）、$HU\le H(K_S+L)\le2$。
$0\le U^2\le (HU)^2/8<1$ より $U^2=0$ で自由な鉛筆である。
$HU=1$ なら $L,U$ の包への直交射影で $H^2\le 2(HL)(HU)/(LU)-L^2(HU)^2/(LU)^2=5$、矛盾。
従って $HU=2$。$H-L-U$ は $H$ に直交して自乗零なので $H\equiv L+U$。
ゆえに入力縮約曲線はすべて $L$ 次数零である。ところが跡の有理切断は
$L$ との交点 $2-2g_L=-2$ を持ち、極は入力縮約曲線に限られる。これは不可能である。

$g_L=1$ の場合も随伴から $K_S+L$ は有効で $L$ に直交する。
最小解消の例外成分で $K_S$ の次数は非負である
（負の自乗と随伴より例外的な滑らかな $(-1)$ 曲線のみが例外となり得る）。
従って $K_S+L$ の有効因子は自己の全成分に非負に交わり、負定値性より零、$K_S=-L$。
定理\ref{thm:cofactor-eighteen-no-smooth-quadric}と同様、
$H+K_S$ の種数二随伴鉛筆から $H=L+U$ で $U$ は自由な有理鉛筆、
跡の零と極は全て $L$ 次数零の曲線に支えられる。

核の特異点の種類を仮定せず、これらの例外曲線の入力像で評価を調べる。
核の任意の一点 $y$ を通る平面系の引戻しの固定曲線部分を $Z_y$ とすると
\[
        (L-Z_y)^2\ge4-\operatorname{mult}_yK\ge1 .
\]
実際二つの一般平面の共通直線が $y$ の外で切る点を数えればよく、重複度四なら錐である。
$LZ_y=0$, $(H-3L/2)^2=-1$ なので $0\le HZ_y\le\sqrt3<2$。
各例外核点の上の一般入力切断は、有限個の孤立逆像を避け、例外曲線に会う場合も
そのただ一つの一般点で横断的に会い、一般核座標の横断位数は一である。
従って核の微分は切断上そこで非零である。該当曲線の入力像は直線なので一般には階数三で、
この入力点での正規化繊維は長さ一である。他の曲線の一般点では既に跡が同型である。
ゆえに一般入力平面で評価は正規化の押出しに等しいが、
そのEuler数は次数六、種数二から五で $M_H$ の四に反する。
\end{proof}

\begin{proof}[有理正規スクロール]
$g_L=0$ では一般切断の制限と $H^1(\mathcal O_S)=0$ から $h^0(L)=6$。
完全系の像は $\mathbb P^5$ 内の次数四の極小次数曲面で、その正規性と
入力した四核座標への有限線形射影から $K$ の正規化である。
極小次数曲面の分類では平面の二次Veroneseまたは有理正規スクロール（錐を含む）。
平面モデルでは $L=2h$, $HL=6$ から逆移動部は $3h$ で自乗九になり、前命題に反する。
$\mathbb F_j$ 上のスクロールは $j=0,2,4$、$L=E+(4+j)F/2$ である。
$j=4$ は完全系の像が錐で、外部中心からの線形射影も錐になる。
$H=aE+bF$ とおくと $H^2=12a-4a^2=8$ より $a=1,2$ となる。

$a=1$ なら $b=4+j/2$、$H$ は豊富で跡因子は $2E+(j-2)F$。
$j=0$ では有効でなく、$j=2$ では $2E$ そのものである。
後者で核射は $E$ を直線へ同型に写す。$E$ の一般点で核の微分は単射である。
実際核母線方向とこの直線の方向が一致すれば、その母線像はこの直線自身となり、一般には起きない。
$E$ の入力像は次数三以下なので階数三の一般点で核は入力から復元される。
その上に他の一般正規化枝があればその枝でも跡は非単元である。
実際階数三で $M$ は可逆かつ評価で $L$ を生成するため、跡因子は導手因子であり、
分離されない枝のいずれでも導手は単元にならない。
台は $E$ だけ、しかも核座標が $E$ 上で単射なので他の枝はない。
よって有限入力射がそこで同型となり、跡の消滅に反する。

$a=2$ では $H=2E+(2+j)F$、跡切断の類は $2F$。
各核母線は直線に同型、各入力母線像は滑らかな円錐曲線で、そこへの射も同型である。
実際次数二の自由系が直線へ二対一なら、同一入力の二点の核座標が異なるため
その入力直線が階数二軌跡に入ってしまう。
$j=2$ のとき $H$ が潰すのは $E$ だけである。$E$ を除いた跡の有効零因子の閉包は
$2F+kE$ で $k\ge0$（母線次数）、$k\le1$（$E$ との交点）。
$k=1$ ならこの閉包は $E+2F$ の滑らかな切断 $T$ そのものである。
実際 $E$ 成分を持たず他の水平成分が必要で、その類だけで全次数を使う。
$LT=3$, $HT=4$、入力像は次数四以下で一般には階数三である。
$T$ から入力像への次数は $HT,LT$ を共に割るから一であり、
跡の台がこの切断だけなので他の正規化枝はない。
その一般点で導手の位数が一となるのは不可能である。一般横断平面で一つの滑らかな正規化枝があり、
平面Gorenstein曲線の導手次数は $2\delta$ である
（有限双対で正規化余長と導手から元の環までの余長が等しい）。従って $k=0$。

残る両 $j$ では、一般に有限な入力開集合での跡の零は二母線の和である。
相異なる二本なら各係数は一であり、一枝だけなら前段落の偶奇に反する。
従って二母線の入力円錐曲線が同じで、階数三での復元からその対応点の核像も一致する。
二重母線ならその円錐曲線上の一般点は一枝だけであり、$f$ の微分の階数は一である。
さもなくば有限射がその一般点で同型となる。また $y$ は入力から復元されるので、
入力射と核射の対の微分もその母線上で階数一以下である。この閉条件は母線全体で成立する。

$j=2$ なら $f(E)$ は一点、核射は $E$ 上の次数一の埋込みである。
相異なる二母線が対応するなら $E$ との各交点は同じ $f(E)$ に写り、円錐曲線のパラメータから一意に対応する。
二点の核像も同じとなり $E$ 上の埋込みに反する。
二重母線なら $E$ との交点で、$E$ の接方向は入力で消えて核では消えず、
母線の接方向は入力で消えない。従って対の微分階数が二となり矛盾する。

$j=0$ なら $S=\mathbb P^1_s\times\mathbb P^1_t$ とし $t$ を母線パラメータとする。
相異なる二母線の円錐曲線表示の対応は $\mathbb P^1_s$ の射影同型で、固定点 $s_0$ を持つ。
切断 $s=s_0$ の入力次数は二で、相異なる二点を同じ点に写すから直線へ二対一となる。
二重母線 $t=t_0$ の場合もこのような切断がある。底方向の微分は円錐曲線の接束に全て入り、
母線方向の微分を使って $\mathbb P^1_s$ の大域ベクトル場になる
（$t_0$ で非零の底接ベクトルを選ぶ）。このベクトル場の零点を $s_0$ とすれば、
切断の入力写像の微分がそこで零で、やはり次数二の切断は直線に二対一である。
その入力直線上の一般点は階数三だが二正規化点を持つので、両枝で導手が非単元でなければならない。
ここで見つけた切断はどちらの母線とも異なり跡因子の外なので矛盾する。
これで四次も全て排除された。
\end{proof}
\path{computations/verify_cofactor_quartic_normalization.py} は線織格子、種数二鉛筆および最小次数曲面上の全数値型を検算する。
